% ID: SEL-D0001-OB-E04
% ESTADO: validada. S002, 2026-11-07. Tres pasadas del mismo autor; ver informe.
\textbf{Datos y teoría.} El tiempo poblacional es Normal y su desviación típica es conocida: $\sigma=8$ minutos. Por ello se emplea un percentil de la Normal estándar para la media, incluso siendo pequeña la muestra; no se sustituye $\sigma$ por una desviación muestral ni se usa una distribución $t$.

La suma de los nueve tiempos es 135 minutos y la media muestral es
\[
\overline x=\frac{10+17+8+27+6+9+32+5+21}{9}=\frac{135}{9}=15\ \text{minutos}.
\]
\textbf{a) Intervalo de confianza del $92\%$.}\par\nopagebreak[4]

La probabilidad de cada cola es $(1-0.92)/2=0.04$. Buscamos $\Phi(z)=0.96$, de donde $z\approx1.750686$ (aproximadamente $1.75$ con tabla). El error de estimación es
\[
E=z\frac\sigma{\sqrt n}=1.750686\ldots\frac8{\sqrt9}
\approx4.668496\ \text{minutos}.
\]
Así,
\[
IC_{92\%}(\mu)=\left(\overline x-E,\overline x+E\right)
\approx\boxed{(10.331504,\ 19.668496)\ \text{minutos}}.
\]
El semiancho es $E$; la amplitud total sería $2E$.

\textbf{b) Tamaño mínimo con confianza del $95.5\%$.}\par\nopagebreak[4]

La confianza cambia, por lo que debemos calcular un nuevo percentil:
\[
\alpha=0.045,\qquad 1-\frac\alpha2=0.9775,
\qquad z=\Phi^{-1}(0.9775)\approx2.004654462.
\]
Con la tabla, $0.9775$ queda entre los valores de $2.00$ y $2.01$; una interpolación da aproximadamente $2.005$. Conservamos más cifras para decidir el mínimo. Sustituir sin más el percentil por 2 modificaría esa decisión.

El error debe ser estrictamente menor de $1.5$ minutos:
\[
z\frac8{\sqrt n}<1.5
\quad\Longleftrightarrow\quad
n>\left(\frac{8z}{1.5}\right)^2
\approx114.307968.
\]
Por tanto,
\[
\boxed{n_{\min}=115\text{ empleados}.}
\]
\textbf{Comprobación del mínimo.} Con el percentil sin redondear,
\[
E(114)\approx1.502025>1.5,\qquad
E(115)\approx1.495480<1.5.
\]
El error disminuye con el tamaño muestral, por lo que 115 es el menor tamaño que cumple la condición con la confianza solicitada.
